%% Renamed from "Background and Related Work" to the canonical A* section name.
\section{Related Work}
\label{sec:related}

Default LLM tutors reward fluent paragraphs. That design assumes the learner
can hold the whole turn in working memory, stay on one topic, and re-enter
after a tab switch. For many students with ADHD, those assumptions fail in
ordinary use. Prior work in three areas makes that failure predictable, yet
none of it closes the gap that ADHD Assist addresses.

\subsection{ADHD as a dimensional, multi-deficit condition}

Start with the condition itself: ADHD is not one ``attention off'' switch.
DSM-5 presentations (inattentive, hyperactive-impulsive, combined) sit on
severity ranges, and executive-function accounts go
deeper~\cite{apa2022dsm}. Barkley treats behavioral inhibition as core, with
secondary damage to working memory, internalized speech, affect and motivation
regulation, and reconstitution~\cite{barkley1997behavioral}. Brown clusters
activation, focus, effort, emotion, memory, and action~\cite{brown2013new}.
Adult evidence treats emotional dysregulation as more than a side
effect~\cite{beheshti2020emotion}. Sonuga-Barke's dual pathway adds delay
aversion and reward-motivation alongside executive
dysfunction~\cite{sonuga2003dual}, so interest and near-term payoff can mask
or amplify the same deficits.

\paragraph{What we take.}
Tutoring has to name \emph{which} deficit a reply shape is meant to repair
(working memory, initiation, impulsivity, set-shifting, affect), not only
``make chat nicer.''
\paragraph{What we cannot claim from taxonomy alone.}
That any particular UI or prompt \emph{works} until it is measured under use.

The closest GenAI bridge is Zhu, Yu, and Luo~\cite{zhu2026scaffolding}. ADHD
university students and experts co-design scaffolding that limits the visible
agenda, prefers progressive tasks, and warns against cognition outsourcing and
clinical overclaim. We take their interaction directions as primary
human-grounded inspiration for the pillars. We do not treat their study as a
measured evaluation of EduAI.

\subsection{Cognitive load and accessibility foundations}

Even if the diagnosis is clear, packaging still matters. Cognitive load theory
separates intrinsic task demand from extraneous packaging and germane
work~\cite{sweller2011cognitive}. Extraneous load rises when answers dump
undifferentiated text. Working memory typically holds only a few chunks at
once~\cite{cowan2010magical}. Accessibility guidance makes that concrete:
W3C~COGA pushes clear purpose, hierarchy, plain language, and predictable
structure, and names AD(H)D among populations that
benefit~\cite{w3c2021coga}. Mayer's segmenting, signaling, and coherence
principles~\cite{mayer2021multimedia}, and UDL's multiple means of
representation~\cite{cast2018udl}, point the same way: signpost, chunk, and do
not overwhelm the first surface.

\paragraph{What we take.}
Length caps, summary-first layout, and one-decision-at-a-time disclosure are
load-management techniques with established warrants. They are not decorative
``AI style.''
\paragraph{What we cannot claim.}
That CLT or COGA alone prove ADHD-specific superiority over neurotypical
learners. Those frameworks are general. Whether Assist helps students with
ADHD \emph{more} than peers is powered human follow-on work, not this paper's
primary result.

\subsection{LLM tutoring architecture and enforcement}

A third literature designs LLM tutors as \emph{pipelines}, not single prompts.

Liu et~al., SocraticLM, introduce a Dean--Teacher--Student pattern: a Dean
judges whether a generated turn meets stated requirements and may revise it
before presentation to the learner~\cite{liu2024socraticlm}. They also score
tutor readability, productive refusal of irrelevant moves, and
validate-and-move. We take the oversight \emph{cadence} and the
readability~/~redirect vocabulary. We do \textbf{not} claim their Socratic
fine-tuning gains or those score numbers as ours. Those metrics stay parked
for later EduAI instrumentation.

Saha, Hase, and Bansal show selective explanation can beat always-explain,
that style exemplars beat long rule lists for compliance, and that a
\emph{misaligned} teacher harms the trajectory~\cite{saha2023can}. This is
perhaps the strongest evidence that prompting alone is fragile. Shani et~al.\ treat a
written \emph{constitution} of tutoring quality over full dialogues and treat
multi-topic teacher turns as violations~\cite{shani2024multiturn}; our
single-focus pillar and Dean constitution borrow that framing. Chevalier
et~al.\ warn that dialogue or style tuning without content checks can break
subject accuracy~\cite{chevalier2024lmscience}, so oversight here must aim at
\textbf{structure}, not inventing facts (key-point style grading motivates our
content-parity discipline). Choudhury and Sodhi (\emph{Better than your
teacher}~/~LEAP) argue for privileged teachers that see full drafts before a
weak agent acts~\cite{choudhury2025better}; we take the
full-draft-before-stream idea only as architectural analogy. Their benchmarks
are not tutoring.

Ma et~al.\ support simulated-learner QA infrastructure~\cite{ma2025students}.
Big Five is not ADHD, so that line is useful for harnesses, not clinical
claims. Heterogeneous-learner teaching theory (e.g., Zhang et~al.\
MINT~\cite{zhang2023nonparametric}) motivates personalisation only at a high
level. It does not justify a user-facing ADHD rule by itself.

\subsection{Gap}
\label{sec:gap}

The components needed for a solution already exist in prior work. Zhu et~al.\
show what ADHD-friendly scaffolding looks like in design workshops. Sweller,
Cowan, and the W3C~COGA guidance explain why short, hierarchical, predictable
replies reduce extraneous load. Saha et~al.\ and Shani et~al.\ show that
teaching style needs selective policy and a written constitution, because
open prompting is fragile. Liu et~al.\ (and LEAP-style full-draft review)
show that a second agent can judge a turn and revise it before the learner
sees it.

What is missing is a single working tutoring control path that combines them.
Two connections are required, and no prior system provides both. The first is
design clarity: assembling the already-known technique-to-deficit warrants
(from ADHD research, load theory, and co-design) into one explicit technique
$\times$ deficit mesh for a tutoring product, and treating that mesh as the
design constitution (\S\ref{sec:design}). The mapping is literature-grounded;
Study~1 does not re-prove every cell as a clinical outcome. The second is
system evaluation: writing that mesh as a runtime policy, adding a verifying
second agent, and measuring structural adherence under multi-turn drift in
three arms (baseline, prompt-only, oversight) while holding model, retrieval,
and tools fixed. That is the appropriate test of whether enforcement beats
prompting when dialogue context pulls the tutor off-scaffold.

The closest prior systems stop one step short. Zhu et~al.\ motivate the
patterns, but they do not run an enforceable Dean on a live tutor under drift
probes. SocraticLM does include a Dean, yet its constitution is Socratic
pedagogy, not ADHD load~/~structure rules, and it does not report our
prompt-vs-oversight ablation. Selective-teaching and RLHF-education papers
motivate policy and audit in general; they do not target ADHD scaffolds or
late-turn structural pass under interruptions.

This paper addresses that gap. We design the mesh, instantiate a
Router / Teacher / Student / Dean control path in EduAI as
control over reply shape, and measure whether oversight beats prompting under
multi-turn drift. Human load and learning stay feasibility-only until we have
powered evidence.
